\documentclass[11pt,a4paper]{article}

\usepackage[T1]{fontenc}
\usepackage[utf8]{inputenc}
\usepackage[italian]{babel}
\usepackage{lmodern}
\usepackage[margin=2.5cm]{geometry}
\usepackage{parskip}   % paragrafi senza rientro, separati da uno spazio verticale
\usepackage{amsmath}
\usepackage{booktabs}
\usepackage[hidelinks]{hyperref}

\title{Steam Recommender\\[4pt]
\large Confronto tra content-based e collaborative filtering con feedback implicito}
\author{Alessandro Casale -- Matricola 578749\\
Sistemi Intelligenti per Internet\\
Laurea magistrale in Ingegneria Informatica, Università Roma Tre}
\date{}

\begin{document}

\maketitle

\begin{center}
Codice: \url{https://github.com/casalea040-cloud/steam-recommender}
\end{center}

\section{Introduzione}

In questo progetto ho realizzato un sistema di raccomandazione di videogiochi per
gli utenti di Steam, con l'obiettivo di confrontare i due approcci visti a
lezione: il content-based e il collaborative filtering. Ho applicato entrambi
agli stessi utenti e giochi e li ho valutati con le stesse metriche,
così da poterne confrontare risultati e limiti.

Ho usato i dati pubblici di Steam raccolti dalla UC San Diego, che contengono tre
informazioni: i giochi presenti nella libreria di ogni utente, con il tempo di
gioco; le caratteristiche dei giochi (generi, tag, modalità di gioco); le
recensioni scritte dagli utenti. Mentre il tempo di gioco, espresso in minuti, è presente per ogni coppia utente-gioco, le recensioni sono molte meno. Per questo ho dovuto ricavare dalle ore di gioco un \emph{indice di
gradimento}, utile per entrambi gli approcci affrontati: un gioco è considerato gradito se l'utente lo ha recensito
positivamente oppure, in assenza di recensione, se ci ha giocato almeno due ore
(Sezione~\ref{sec:gradimento}). 

\subsection*{I due approcci}

Nel \textbf{content-based} ogni gioco è rappresentato da un vettore costruito dai
suoi generi, tag e modalità di gioco. Il profilo di un utente è la media dei
vettori dei giochi che ha gradito; a ogni gioco che non possiede assegno come
punteggio la similarità coseno con il suo profilo, e gli raccomando i giochi con
il punteggio più alto. Questo approccio usa solo le caratteristiche dei giochi e
i dati dell'utente stesso.

Nel \textbf{collaborative filtering} le caratteristiche dei giochi non vengono
usate: conta solo chi ha gradito cosa, cioè una matrice con gli utenti sulle righe e i
giochi sulle colonne (Sezione~\ref{sec:matrici}). Ho considerato due versioni. Nella versione
\emph{user-based} confronto le righe della matrice per trovare gli utenti con
gusti più simili, e raccomando i giochi che loro hanno gradito. Nella versione
\emph{item-based} confronto le colonne, considerando simili due giochi graditi
dalle stesse persone, e raccomando i giochi più simili a quelli che l'utente ha
già gradito. In entrambi i casi uso la similarità coseno.

\section{Dataset}

Ho usato lo Steam Video Game and Bundle Data della UC San
Diego~\cite{dataset}, e in particolare tre file:
\begin{itemize}
  \item \texttt{australian\_users\_items}: per ogni utente, i giochi posseduti
        con il tempo di gioco in minuti;
  \item \texttt{steam\_games}: per ogni gioco, titolo, generi, tag, specs,
        prezzo e data di uscita;
  \item \texttt{australian\_user\_reviews}: le recensioni, con l'indicazione
        consigliato sì/no.
\end{itemize}
I file sono collegati dall'identificativo del gioco. La
Tabella~\ref{tab:dataset} riassume i dati prima della pulizia.

\begin{table}[h]
\centering
\begin{tabular}{lr}
\toprule
Coppie utente--gioco nelle librerie & 5.153.209 \\
Utenti con almeno un gioco & 70.912 \\
Giochi nelle librerie & 10.978 \\
\quad di cui con metadati & 9.192 \\
Coppie utente--gioco con 0 minuti di gioco & 36,2\% \\
Recensioni (da 25.458 utenti) & 59.305 \\
\quad di cui positive & 88,5\% \\
\bottomrule
\end{tabular}
\caption{Il dataset prima della pulizia.}
\label{tab:dataset}
\end{table}

Più di un gioco posseduto su tre non è mai stato avviato, probabilmente perché
acquistato in saldo o all'interno di un bundle: il possesso da solo non indica
gradimento. Per descrivere i giochi ho scelto come fonte principale i tag,
assenti solo nello 0,5\% dei casi contro il 10,2\% dei generi; il dataset non
contiene descrizioni testuali.

\section{Acquisizione e memorizzazione dei dati}

Ho sviluppato il progetto in Python, usando pandas per le tabelle, SciPy e
scikit-learn per le matrici e le similarità, PyArrow per il formato Parquet e
Streamlit per la web app.

Ho scritto uno script, \texttt{download\_data.py}, che scarica i file in
\texttt{data/raw}. I file non sono JSON standard (ogni riga è un dizionario
Python), quindi nel modulo \texttt{load.py} li leggo riga per riga con
\texttt{ast.literal\_eval} e li trasformo in tabelle pandas in formato lungo: una
riga per coppia utente--gioco, per recensione e per gioco. Durante la lettura
uniformo gli identificativi, rendo generi, tag e specs sempre liste e converto il
prezzo in numero (0 per i giochi gratuiti).

La lettura dei file originali richiede alcuni minuti, quindi salvo le tabelle in
formato Parquet in \texttt{data/interim} e nelle esecuzioni successive le rileggo
da lì in pochi secondi. Le tabelle pulite le salvo in \texttt{data/processed}.
Non modifico mai i file originali, così l'intero processo si può rieseguire dagli
script. Ho raccolto i percorsi in un unico modulo di configurazione e ho
pubblicato il codice su GitHub, lasciando fuori dal repository i dati.

\section{Pulizia dei dati}

Per la pulizia ho fatto queste considerazioni:
\begin{itemize}
  \item elimino le coppie utente--gioco e le recensioni duplicate;
  \item escludo i giochi senza metadati (1.786 dei 10.978 presenti nelle
        librerie), così che i due approcci lavorino sugli stessi giochi;
  \item dopo aver calcolato l'indice di gradimento
        (Sezione~\ref{sec:gradimento}), escludo gli utenti con meno di 5 giochi
        graditi e i giochi graditi da meno di 5 utenti.
\end{itemize}

Le due soglie si influenzano a vicenda: togliere un gioco può far scendere un
utente sotto i 5 giochi graditi, e togliere un utente può far scendere un gioco
sotto i 5 utenti. Per questo ripeto i due filtri finché un passaggio non toglie
più nulla. La Tabella~\ref{tab:pulizia} mostra l'effetto di ogni passo.

\begin{table}[h]
\centering
\begin{tabular}{lrrr}
\toprule
Passo & Utenti & Giochi & Coppie utente--gioco \\
\midrule
Librerie originali             & 70.912 & 10.978 & 5.153.209 \\
Eliminazione dei duplicati     & 70.912 & 10.978 & 5.094.082 \\
Solo giochi con metadati       & 68.712 &  9.192 & 4.244.812 \\
Filtri, primo passaggio        & 53.659 &  5.230 & 4.020.111 \\
Filtri, secondo passaggio      & 53.653 &  5.229 & 4.019.999 \\
Filtri, terzo passaggio        & 53.653 &  5.229 & 4.019.999 \\
\bottomrule
\end{tabular}
\caption{Utenti, giochi e coppie utente--gioco dopo ogni passo della pulizia.}
\label{tab:pulizia}
\end{table}

Quasi tutta la riduzione avviene al primo passaggio dei filtri, che toglie circa
15.000 utenti e 4.000 giochi: molti utenti hanno pochi giochi graditi, e molti
giochi sono apprezzati da pochissime persone. Il secondo passaggio toglie solo 6
utenti e un gioco, effetto a catena del primo, e il terzo non cambia più nulla,
quindi il ciclo si ferma. Alla fine restano 53.653 utenti, 5.229 giochi e
4.019.999 coppie utente--gioco posseduto, di cui 1.676.503 gradite.

\section{Indice di gradimento}
\label{sec:gradimento}

Poiché possedere un gioco non significa averlo apprezzato, ho definito una regola
che stabilisce quando un gioco è gradito. Indicando con $t(u,i)$ il tempo di
gioco in minuti dell'utente $u$ sul gioco $i$, l'indice che ho adottato è:
\[
g(u,i) =
\begin{cases}
1 & \text{se } u \text{ ha recensito positivamente } i,\\
0 & \text{se } u \text{ ha recensito negativamente } i,\\
1 & \text{se non c'è recensione e } t(u,i) \geq 120,\\
0 & \text{altrimenti.}
\end{cases}
\]
Ho dato la precedenza alla recensione, quando presente: essendo un giudizio
espresso direttamente dall'utente, deve essere più importante di qualsiasi altra valutazione estrapolata dai dati.

Ho fissato la soglia a due ore (espresse in minuti) riprendendo direttamente la politica di rimborso di Steam, che
consente di restituire un gioco entro due settimane dall'acquisto se è stato
giocato meno di due ore~\cite{refund}. Chi supera questa soglia ha continuato a
giocare oltre il periodo in cui poteva ancora chiedere il rimborso, e lo
considero un segnale di gradimento. In questo modo i giochi mai avviati risultano
non graditi, e la quota di giocatori che gradiscono un gioco varia da gioco a
gioco.

Inizialmente avevo considerato una soglia relativa: la mediana delle ore dei
giocatori di ciascun gioco, che rendeva  confrontabili giochi brevi e lunghi. Con la mediana l'idea era fissare valore gradimento 1 se l'utente aveva un numero di ore maggiore del valore della mediana, 0 altrimenti. L'ho
scartata perché, per costruzione, classifica come non gradito ciascun gioco a esattamente metà dei giocatori, indipendentemente da quanto il gioco sia apprezzato in realtà. La soglia
fissa ha il limite opposto: due ore pesano allo stesso modo per un gioco breve e
per uno molto lungo.

\section{Modelli}
\label{sec:modelli}

\subsection{La matrice utenti $\times$ giochi}
\label{sec:matrici}

Tutti i modelli lavorano su matrici con gli utenti sulle righe e i giochi sulle
colonne. Le costruisco a partire dalla tabella pulita, in cui ogni riga è una
coppia utente--gioco con i minuti di gioco e l'indice di gradimento, e ne ricavo
due:
\begin{itemize}
  \item la matrice di gradimento $G$, in cui ogni cella contiene l'indice di
        gradimento $g(u,i)$ dell'utente per quel gioco, e vale 0 per i giochi
        che l'utente non possiede (Tabella~\ref{tab:matrice});
  \item la matrice di possesso $O$, che vale 1 per ogni gioco presente nella
        libreria dell'utente, gradito o no.
\end{itemize}

\begin{table}[h]
\centering
\begin{tabular}{lccc}
\toprule
 & Gioco 1 & Gioco 2 & Gioco 3 \\
\midrule
Utente A & 1 & 0 & 0 \\
Utente B & 1 & 1 & 0 \\
Utente C & 0 & 1 & 1 \\
\bottomrule
\end{tabular}
\caption{Esempio di matrice di gradimento $G$: 1 = gioco gradito, 0 =
gioco non gradito o non posseduto.}
\label{tab:matrice}
\end{table}

Uso $G$ per calcolare similarità e punteggi; $O$ mi serve per non raccomandare a
un utente giochi che possiede già, anche quelli che non ha gradito. Con decine
di migliaia di utenti e migliaia di giochi le matrici sono quasi tutte vuote,
perché ogni utente possiede una piccola parte del catalogo; per questo le
memorizzo come matrici sparse, salvando solo le celle diverse da zero.

\subsection{Collaborative filtering user-based}

L'idea del collaborative filtering user-based è che utenti che in passato hanno
gradito gli stessi giochi probabilmente continueranno a gradire gli stessi. Per raccomandare giochi a un utente cerco quindi gli utenti con i
gusti più simili ai suoi, i suoi \emph{vicini}, e gli propongo i giochi che loro
hanno gradito e lui non possiede.

\paragraph{Similarità tra utenti.} Rappresento ogni utente $u$ con la sua riga
$\mathbf{g}_u$ di $G$ e misuro la somiglianza tra due utenti con la similarità
coseno. Poiché le righe contengono solo 0 e 1, indicando con $I_u$ l'insieme dei
giochi graditi da $u$, la formula ha un significato semplice:
\[
\mathrm{sim}(u,v) = \frac{\mathbf{g}_u \cdot \mathbf{g}_v}
{\lVert\mathbf{g}_u\rVert\,\lVert\mathbf{g}_v\rVert}
= \frac{|I_u \cap I_v|}{\sqrt{|I_u|\,|I_v|}}.
\]
Il numeratore conta i giochi graditi da entrambi; il denominatore evita che un
utente con una libreria enorme risulti simile a tutti solo perché ha gradito
molti giochi. La similarità vale 0 se i due utenti non hanno giochi graditi in
comune e 1 se hanno gradito esattamente gli stessi.

\paragraph{Vicini e punteggi.} Per ogni utente $u$ prendo i $k$ utenti con
similarità più alta, che indico con $N_k(u)$; per ora ho fissato $k = 50$. A
ogni gioco $i$ assegno il punteggio
\[
s(u,i) = \frac{\sum_{v \in N_k(u)} \mathrm{sim}(u,v)\, g(v,i)}
{\sum_{v \in N_k(u)} \mathrm{sim}(u,v)},
\]
cioè la frazione dei vicini che hanno gradito $i$, in cui ogni vicino pesa
quanto è simile a $u$. Il punteggio è compreso tra 0 e 1: vale 1 se tutti i
vicini hanno gradito il gioco. Infine escludo i giochi che $u$ possiede già,
usando $O$, e raccomando gli $n$ giochi con il punteggio più alto.

Molti giochi hanno punteggio 0, perché nessun vicino li ha graditi. Per
decidere l'ordine tra giochi con lo stesso punteggio aggiungo un contributo
molto piccolo proporzionale alla popolarità del gioco, cioè al numero di utenti
che lo hanno gradito: non cambia l'ordine tra giochi con punteggi diversi, ma
tra due giochi a pari merito fa preferire quello più apprezzato in generale.
Questo copre anche il caso, raro, di un utente senza nessun vicino con
similarità maggiore di 0.

\paragraph{Calcolo a blocchi.} Per trovare i vicini devo confrontare ogni
utente con tutti gli altri. Calcolare in una volta sola tutte le similarità
vorrebbe dire costruire una matrice $53.653 \times 53.653$, quasi 2,9 miliardi
di valori e più di 11~GB di memoria. Per questo elaboro gli utenti a blocchi di
500: per ogni blocco calcolo le similarità con tutti gli utenti, ricavo i vicini
e le raccomandazioni, e tengo in memoria solo queste ultime prima di passare al
blocco successivo. Prima del calcolo divido ogni riga di $G$ per la sua norma, ottenendo $G_n$,
così la similarità coseno diventa un semplice prodotto scalare tra righe.

\subsection{Baseline di popolarità}

Come termine di paragone uso un modello che non personalizza: raccomanda a
tutti gli utenti i giochi graditi dal maggior numero di persone, escludendo
quelli che ognuno possiede già. Su Steam pochi giochi molto famosi compaiono
nelle librerie di quasi tutti, quindi è una baseline difficile da battere: un
modello personalizzato ha senso solo se riesce a fare meglio.

\section{Valutazione e web app}

Per valutare gli approcci, per ogni utente nascondo una parte dei giochi graditi
e misuro quanti di questi il sistema riesce a ritrovare tra le prime $n$
raccomandazioni, con precision e recall. Come termine di paragone uso la
baseline di popolarità descritta sopra. Oltre al confronto complessivo,
analizzerò il comportamento sui giochi con pochi giocatori, dove il
collaborative filtering ha poche informazioni, e la varietà dei giochi
raccomandati.

Infine affiancherò alla valutazione una piccola web app che, scelto un utente,
mostra una accanto all'altra le raccomandazioni dei due approcci, ciascuna con
una breve spiegazione del perché il gioco è stato proposto.

\textbf{[Da continuare]}

% ---------------------------------------------------------------------------
% Sezioni da completare:
% \subsection{...}          content-based (tag binari e TF-IDF) e CF item-based
% \section{Valutazione}     split training/test sui giochi graditi, precision e
%                           recall sulle prime n, cold start, varietà;
%                           codice di evaluate.py per il confronto dei modelli
% \section{Web app}         Streamlit, raccomandazioni dei due approcci affiancate
% \section{Conclusioni}
% ---------------------------------------------------------------------------

\begin{thebibliography}{9}

\bibitem{dataset}
J.~McAuley. \emph{Recommender Systems and Personalization Datasets}. UC San
Diego. \url{https://cseweb.ucsd.edu/~jmcauley/datasets.html}

\bibitem{refund}
Valve. \emph{Steam Refunds}. \url{https://store.steampowered.com/steam_refunds}

\end{thebibliography}

\end{document}
